\chapter{切空间与微分形式}
\label{chap:tangent-differential-forms-solutions}

\begin{solution}{tangent-coordinate-basis-change}
设 \(y=y(x)\) 是两组局部坐标之间的变换。把
\(\partial/\partial x^i|_p\) 看成导子。对任意局部光滑函数 \(f\)，链式
法则给出
\[
 \frac{\partial f}{\partial x^i}(p)
 =\sum_j\frac{\partial y^j}{\partial x^i}(p)
       \frac{\partial f}{\partial y^j}(p).
\]
因为该式对每个 \(f\) 成立，所以
\[
 \left.\frac{\partial}{\partial x^i}\right|_p
 =\sum_j\frac{\partial y^j}{\partial x^i}(p)
       \left.\frac{\partial}{\partial y^j}\right|_p.
\]

若
\[
 v=\sum_i v_x^i\left.\frac{\partial}{\partial x^i}\right|_p
  =\sum_j v_y^j\left.\frac{\partial}{\partial y^j}\right|_p,
\]
代入基变换公式并比较 \(\partial/\partial y^j|_p\) 的系数，得到
\[
 v_y^j
 =\sum_i\frac{\partial y^j}{\partial x^i}(p)v_x^i.
\]
换言之，分量向量按照坐标变换的 Jacobian 变换；若把新的坐标基用旧基表示，
则出现逆 Jacobian。两种变化互相抵消，所得导子 \(v\) 与坐标无关。
\end{solution}

\begin{solution}{tangent-hadamard-derivation}
把坐标平移后设 \(x(p)=0\)，并把坐标域缩小成关于径向线段星形的邻域。
记 \(\widetilde f=f\circ x^{-1}\)。对固定的 \(u\)，一元微积分基本定理
给出
\begin{align*}
 \widetilde f(u)-\widetilde f(0)
 &=\int_0^1\frac{d}{dt}\widetilde f(tu)\,dt\\
 &=\sum_i u^i\int_0^1\partial_i\widetilde f(tu)\,dt.
\end{align*}
令
\[
 g_i(u)=\int_0^1\partial_i\widetilde f(tu)\,dt.
\]
积分号下逐阶求导表明 \(g_i\) 光滑，并且
\[
 g_i(0)=\partial_i\widetilde f(0).
\]

由 \(D1=D(1\cdot1)=2D1\) 得 \(D1=0\)，所以导子在常数函数上为零。
把上面的分解写成
\[
 f-f(p)=\sum_i x^i g_i
\]
并使用 Leibniz 规则，注意 \(x^i(p)=0\)，便有
\[
 Df
 =\sum_i\bigl(D(x^i)g_i(p)+x^i(p)Dg_i\bigr)
 =\sum_iD(x^i)\partial_i\widetilde f(0).
\]
因此 \(D\) 完全由
\((D(x^1),\ldots,D(x^n))\) 决定。这也证明了正文坐标表示的唯一性。
\end{solution}

\begin{solution}{tangent-map-curve-chain-rule}
在曲线表示下，
\[
 d(G\circ F)_p[\gamma]
 =[G\circ F\circ\gamma]
 =dG_{F(p)}[F\circ\gamma]
 =dG_{F(p)}dF_p[\gamma].
\]

在导子表示下，对任意 \(h\in C^\infty_{G(F(p))}(P)\)，有
\begin{align*}
 \bigl[dG_{F(p)}(dF_pv)\bigr]h
 &=(dF_pv)(h\circ G)\\
 &=v(h\circ G\circ F)\\
 &=\bigl[d(G\circ F)_pv\bigr]h.
\end{align*}
两个导子对所有函数芽的作用相同，因而相等。两种证明分别把通常链式法则
表述为复合曲线的速度公式和复合函数的方向导数公式。
\end{solution}

\begin{solution}{tangent-regular-level-kernel}
令 \(\iota:S\hookrightarrow M\) 为包含映射。若
\(v=[\gamma]\in T_pS\)，则 \(F\circ\gamma\equiv c\)，从而
\[
 dF_p(d\iota_pv)=0.
\]
在用 \(d\iota_p\) 把 \(T_pS\) 视为 \(T_pM\) 的子空间后，这给出
\[
 T_pS\subseteq\ker dF_p.
\]

由于 \(c\) 是正则值，正则水平集定理给出 \(p\) 附近的坐标
\((u^1,\ldots,u^{n-m},w^1,\ldots,w^m)\)，使得 \(F\) 在目标的相应
坐标中写成
\[
 (u,w)\longmapsto w,
\]
而 \(S\) 由 \(w=0\) 描述。于是 \(T_pS\) 由
\(\partial/\partial u^1,\ldots,\partial/\partial u^{n-m}\) 张成；
\(dF_p\) 的核也恰由这些向量张成。因此反向包含成立，故
\[
 T_pS=\ker dF_p.
\]
\end{solution}

\begin{solution}{tangent-product-splitting}
令 \(\pi_M,\pi_N\) 为两个投影，定义
\[
 \Psi:T_{(p,q)}(M\times N)\longrightarrow T_pM\oplus T_qN,
 \qquad
 \Psi(v)=(d\pi_Mv,d\pi_Nv).
\]
在积坐标 \((x,y)\) 中，\(\Psi\) 把坐标分量
\[
 (v^1,\ldots,v^m,w^1,\ldots,w^n)
\]
分成前后两块，所以它是线性同构。由于定义只使用两个典范投影，这个同构
不依赖坐标。

若 \(F:M\to M'\)、\(G:N\to N'\)，则
\[
 \pi_{M'}\circ(F\times G)=F\circ\pi_M,\qquad
 \pi_{N'}\circ(F\times G)=G\circ\pi_N.
\]
对这两个等式求切映射，并使用链式法则，得到
\[
 d(F\times G)_{(p,q)}(v,w)=(dF_pv,dG_qw).
\]
这正是所要求的自然相容性。
\end{solution}

\begin{solution}{lie-bracket-coordinate-formula}
写
\[
 X=\sum_iX^i\partial_i,\qquad
 Y=\sum_jY^j\partial_j.
\]
对光滑函数 \(f\) 展开：
\begin{align*}
 X(Yf)
 &=\sum_{i,j}X^i(\partial_iY^j)\partial_jf
   +\sum_{i,j}X^iY^j\partial_i\partial_jf,\\
 Y(Xf)
 &=\sum_{i,j}Y^i(\partial_iX^j)\partial_jf
   +\sum_{i,j}Y^iX^j\partial_i\partial_jf.
\end{align*}
第二行的二阶项交换求和指标 \(i,j\) 后，与第一行的二阶项相同，因为
\(\partial_i\partial_jf=\partial_j\partial_if\)。相减后只剩
\[
 [X,Y]f
 =\sum_j\left(\sum_i
   \bigl(X^i\partial_iY^j-Y^i\partial_iX^j\bigr)\right)\partial_jf.
\]
因此 \([X,Y]\) 是一阶导子，并具有正文所列的坐标表达式。
\end{solution}

\begin{solution}{lie-bracket-jacobi-expansion}
用并置表示算子的复合。第一项展开为
\[
 [X,[Y,Z]]=XYZ-XZY-YZX+ZYX.
\]
循环地，
\begin{align*}
 [Y,[Z,X]]&=YZX-YXZ-ZXY+XZY,\\
 [Z,[X,Y]]&=ZXY-ZYX-XYZ+YXZ.
\end{align*}
三式相加时，六种三重算子复合各出现一次正号和一次负号，故总和为零。由
前一题可知向量场的交换子仍是向量场，所以这个算子恒等式就是 Lie 括号的
Jacobi 恒等式。
\end{solution}

\begin{solution}{lie-bracket-diffeomorphism-naturality}
记
\[
 \widetilde X=F_*X,\qquad \widetilde Y=F_*Y.
\]
对任意 \(h\in C^\infty(N)\)，推前的定义给出
\[
 (\widetilde Xh)\circ F=X(h\circ F),\qquad
 (\widetilde Yh)\circ F=Y(h\circ F).
\]
于是
\begin{align*}
 ([\widetilde X,\widetilde Y]h)\circ F
 &=X\bigl((\widetilde Yh)\circ F\bigr)
   -Y\bigl((\widetilde Xh)\circ F\bigr)\\
 &=X(Y(h\circ F))-Y(X(h\circ F))\\
 &=[X,Y](h\circ F).
\end{align*}
右端等于 \((F_*[X,Y]h)\circ F\)。微分同胚 \(F\) 是满射，因此两个
向量场对所有光滑函数的作用相同，故
\[
 F_*[X,Y]=[F_*X,F_*Y].
\]
\end{solution}

\begin{solution}{cotangent-chain-rule}
对任意 \(v\in T_pM\)，
\[
 d(g\circ F)_p(v)
 =v(g\circ F)
 =dg_{F(p)}(dF_pv)
 =[F^*(dg)]_p(v).
\]
所以
\[
 d(g\circ F)=F^*(dg).
\]
在局部坐标中，若把微分写成行向量，则该式成为
\[
 D(g\circ F)(p)=Dg(F(p))\,DF(p),
\]
即 Jacobian 的链式法则。若把余切向量写成列向量，则相应公式中出现
Jacobian 的转置。
\end{solution}

\begin{solution}{tensor-contraction-invariance}
一个 \((1,1)\) 型张量可以看作线性映射
\[
 A:T_pM\longrightarrow T_pM.
\]
设它在两组坐标基下的矩阵分别为 \(A_x,A_y\)，而基变换矩阵为 \(J\)。
那么
\[
 A_y=J A_xJ^{-1}.
\]
收缩一个上指标与一个下指标，在 \(x\) 坐标中得到
\[
 A^i{}_i=\operatorname{tr}A_x.
\]
利用迹的循环性，
\[
 \operatorname{tr}A_y
 =\operatorname{tr}(JA_xJ^{-1})
 =\operatorname{tr}(A_xJ^{-1}J)
 =\operatorname{tr}A_x.
\]
故该收缩不依赖坐标，并且正是线性算子 \(A\) 的迹。
\end{solution}

\begin{solution}{forms-pullback-wedge-functorial}
设 \(\alpha\) 是 \(k\)-形式，\(\beta\) 是 \(\ell\)-形式。楔积由张量积的
交错化定义。先把每个输入向量经 \(dF\) 送入目标切空间再交错，与分别计算
\(F^*\alpha,F^*\beta\) 后再交错相同。因此
\[
 F^*(\alpha\wedge\beta)=F^*\alpha\wedge F^*\beta.
\]

对复合映射，切映射的链式法则给
\[
 d(G\circ F)_p=dG_{F(p)}\circ dF_p.
\]
逐个代入形式的定义，有
\begin{align*}
 ((G\circ F)^*\omega)_p(v_1,\ldots,v_k)
 &=\omega_{G(F(p))}
   (dG_{F(p)}dF_pv_1,\ldots,dG_{F(p)}dF_pv_k)\\
 &=[F^*(G^*\omega)]_p(v_1,\ldots,v_k).
\end{align*}
故
\[
 (G\circ F)^*=F^*\circ G^*.
\]
这里复合次序的反转来自函数与形式的拉回方向同映射方向相反。
\end{solution}

\begin{solution}{forms-exterior-coordinate-free-formula}
先验证题中右端关于每个 \(X_i\) 都是
\(C^\infty(M)\)-线性的。以 \(X_0\) 为例，把它换成 \(aX_0\) 时，第一
个求和中 \(X_j(a)\) 产生的附加项，与
\[
 [aX_0,X_j]=a[X_0,X_j]-X_j(a)X_0
\]
在第二个求和中产生的附加项恰好相消。其余位置相同。因此右端在一点的取值
只依赖 \(X_0(p),\ldots,X_k(p)\)，确实定义了一个 \((k+1)\)-形式。

固定 \(p\)，选取局部坐标，并把每个 \(X_i(p)\) 延拓成常系数坐标向量场
\(\widetilde X_i\)。由刚才证明的张量性，可以用这些延拓计算题中等式的
两边。常系数坐标向量场彼此交换，第二个求和为零，公式化为
\[
 d\omega(\widetilde X_0,\ldots,\widetilde X_k)
 =\sum_i(-1)^i\widetilde X_i
   \omega(\widetilde X_0,\ldots,\widehat{\widetilde X_i},
          \ldots,\widetilde X_k).
\]
对坐标形式
\[
 \omega=\sum_I a_I\,dx^I
\]
直接展开
\[
 d\omega=\sum_I da_I\wedge dx^I
\]
便得到上式。两边都是张量，所以等式对原来的任意向量场成立。
\end{solution}

\begin{solution}{forms-cartan-one-form-direct}
设 \(\alpha\) 是一次形式，\(Y\) 是任意向量场。由函数外微分以及二形式
外微分的公式，
\begin{align*}
 [d(\iota_X\alpha)+\iota_Xd\alpha](Y)
 &=Y(\alpha(X))+d\alpha(X,Y)\\
 &=Y(\alpha(X))
   +X(\alpha(Y))-Y(\alpha(X))-\alpha([X,Y])\\
 &=X(\alpha(Y))-\alpha([X,Y])\\
 &=(\mathcal L_X\alpha)(Y).
\end{align*}
两边对所有 \(Y\) 的取值相同，故
\[
 \mathcal L_X\alpha=d(\iota_X\alpha)+\iota_Xd\alpha.
\]
\end{solution}

\begin{solution}{forms-flow-lie-derivative}
在题设的共同坐标邻域及时间区间内，局部流联合光滑并满足
\[
 \frac{\partial}{\partial t}\Phi_t(p)=X_{\Phi_t(p)},
 \qquad \Phi_0=\operatorname{id}.
\]
因此，对函数 \(f\)，
\[
 \left.\frac{d}{dt}\right|_{t=0}(\Phi_t^*f)(p)
 =\left.\frac{d}{dt}\right|_{t=0}f(\Phi_t(p))
 =Xf(p).
\]
拉回与外微分可交换；联合光滑性保证坐标中的时间与空间偏导可交换，故
\[
 \left.\frac{d}{dt}\right|_{t=0}\Phi_t^*(df)
 =\left.\frac{d}{dt}\right|_{t=0}d(\Phi_t^*f)
 =d(Xf).
\]

任意微分形式在坐标邻域内都是
\[
 f_0\,df_1\wedge\cdots\wedge df_k
\]
型形式的有限和；这里可取 \(f_j\) 为坐标函数并把系数并入 \(f_0\)。对
拉回后的乘积与楔积求 \(t\) 导数，使用前面两个公式，可知
\[
 A\omega:=\left.\frac{d}{dt}\right|_{t=0}\Phi_t^*\omega
\]
是零次导子，并且在函数及函数的微分上与 \(\mathcal L_X\) 相同。因此
\[
 \left.\frac{d}{dt}\right|_{t=0}\Phi_t^*\omega
 =\mathcal L_X\omega.
\]
再应用 Cartan 公式，得到
\[
 \left.\frac{d}{dt}\right|_{t=0}\Phi_t^*\omega
 =d(\iota_X\omega)+\iota_Xd\omega.
\]
\end{solution}
